\section{Conclusion and Future Work}
\label{sec:conclusion}

In this paper, we presented a rigorous investigation into the post-training quantization (PTQ) robustness of test-time adaptive model merging. We identified and characterized a fundamental, overlooked vulnerability in existing unregularized adaptive merging frameworks: \textbf{Quantization-Operator Overfitting}. While learning layer-wise blending coefficients via test-time adaptation yields high multi-task accuracy in floating-point (FP32) space, it forces the network parameters into sharp local minima that are highly fragile under post-training quantization, triggering catastrophic collapse.
To analyze and mitigate this bottleneck, we presented \textbf{CR-PolySACM (Clipping-Regularized Sharpness-Aware Subspace Model Merging)}, a unified framework that combines global structural constraints with local landscape flatness optimization. First, we identified a fundamental \emph{task-vector norm scale pathology} in standard sharpness-aware optimization, where a 50-fold discrepancy in layer-wise task-vector norms renders standard unnormalized regularizers blind to highly sensitive, low-norm layers such as final layer normalization, causing optimization instability or collapse. Second, we proposed \textbf{Clipping-Regularized SACM (CR-SACM)}, which clips task-vector norms to a robust minimum floor to successfully balance scales without triggering gradient explosion.

Our proposed CR-PolySACM framework integrates CR-SACM with a low-degree depth-dependent polynomial subspace constraint, successfully mapping the test-time adaptation trajectory into flat, robust regions of the landscape. Our comprehensive empirical evaluation across six diverse quantization schemas demonstrates that **CR-PolySACM sets a new state of the art in robust model merging, outperforming the previous state-of-the-art PolyMerge baseline by nearly +1.0\% under severe INT4 quantization (19.07\% vs 18.10\%).** Furthermore, our upgraded HessMerge baseline consistently and significantly outperforms standard AdaMerging across all target precisions (+1.36\% in FP32), demonstrating the critical importance of correcting scale-blindness.

There are several exciting directions for future research. First, we plan to scale these insights to larger architectures such as ViT-Base/Large, vision-language models, and large language models (LLMs), where post-deployment quantization is critical. We hypothesize that these larger, deeper models will exhibit even more pronounced task-vector scale pathologies due to the increased diversity of layer functions (e.g., query-key-value projections, feed-forward networks, and multiple layer normalization blocks) and highly disparate parameter norms. In such regimes, the optimal clipping threshold $\beta$ may need to scale inversely with the average layer-wise task-vector norm, or could be dynamically set as a lower percentile (e.g., the 10th percentile) of the empirical task-vector norm distribution to ensure robust and automated scale balancing across thousands of layers without manual tuning. Specifically, due to the heavy-tailed or highly skewed nature of task-vector norm distributions in LLMs (where a few specialized layers, such as key-value projections or output heads, dominate in norm while the vast majority of intermediate blocks have much smaller, near-identical norms), a fixed clipping threshold $\beta$ is fragile. Under a heavy-tailed distribution, a percentile-based blueprint dynamically partitions the layers: by setting $\beta$ as the 10th percentile of the empirical task-vector norm distribution, we guarantee that only the bottom 10\% of extremely low-norm, highly sensitive layers (the "singular layers" that would otherwise cause gradient explosion or remain completely unseen) are clipped. This automatically isolates and scale-balances the singular tail of the distribution while leaving the mainstream intermediate blocks unperturbed, ensuring seamless generalization and numerical stability regardless of the model scale or depth. We provide a preliminary empirical proof-of-concept evaluation of this percentile-based blueprint on our Vision Transformer backbone in Appendix~\ref{sec:appendix}, where dynamically setting $\beta$ to the 10th percentile achieves $19.05\%$ accuracy under INT4, matching the optimal manually swept threshold. Second, we aim to design more sophisticated scale-balanced optimizers to further refine edge robustness. Finally, we hope to explore methods that directly address the out-of-subspace noise complement $\delta_{\perp}$ during the test-time adaptation phase, such as post-quantization weight reconstruction or scaling-factor calibration, to directly mitigate out-of-subspace noise after the merging coefficients have been optimized.
